Trace retention is an incomplete predictor of observability cost: a sampler changes where work is avoided, how spans are grouped for export, and which evidence remains available. We study these effects in OpenTelemetry Collector Contrib v0.136.0 on one shared node with loopback transport. Five randomized blocks cross sampler placement with export path at 40,000 offered spans/s. Native uniform sampling at nominal 10\% retention increases Collector CPU by \PlacementFileUniformIncrease{}\% with JSON-only export and reduces it by \PlacementDualUniformSaving{}\% with JSON plus Jaeger. A pre-ingress gate retains identical trace IDs but avoids ingestion and excludes selection work from the Collector endpoint. Retain-all controls show that stateful release changes batching and CPU without discarding traces. A short-timeout CPU reduction at 250 spans/s disappears at 5,000 spans/s, where tail increases CPU at both tested timeouts. A separate SDK comparison, load sweep, and delivery probes distinguish application work, Collector resources, and complete evidence delivery. An illustrative localization task uses measured HTTP timings and ideal offline sampling without an SDK or Collector. Its results show how window size and selection-dependent reference evidence affect a fixed median-change scorer, conditional on the observed corpus. The study supports evaluating sampling at explicit component boundaries, measuring batching and serialization alongside trace counts, and defining the diagnostic evidence objective before selecting a rate. The research archive preserves frozen protocols, raw observations, failed attempts, and executable checks.
